\section{Introduction}
\label{a:sec:intro}

For the binary Galton--Watson process with division probability
$p\in[0,\tfrac12)$, the survival probability satisfies
$S_n\sim A(p)(2p)^n$, so the limiting mean population size conditional on
non-extinction is $1/A(p)$. \ChM{} introduced this constant, derived the
conditional mean and variance it controls, and recorded its continuous-time
counterpart $\Ac(p)=(1-2p)/(1-p)$, which is elementary. The discrete constant is
not elementary, and this chapter is devoted to everything that can nevertheless be
said about it exactly.

The results of this chapter are original to this thesis. They comprise an exact
infinite-product representation of $A(p)$ (\cref{a:sec:product}); an exact series
with two-sided bounds that pin $A(p)$ between $\tfrac12\Ac(p)$ and
$2\varepsilon/(1+2\varepsilon)$ (\cref{a:sec:bounds}); a rigorous near-critical
asymptotic with remainder, $A(p)=2\varepsilon\lrb{1+O(\varepsilon\ln(1/\varepsilon))}$
for $\varepsilon=1-2p$ (\cref{a:sec:asymptotic}); the resolution of the
discrepancy between the discrete and continuous constants
(\cref{a:sec:compare}); an account of symbolic-regression searches for an
elementary formula, and of why they failed (\cref{a:sec:GA}); the identification
of $A(p)$ with twice the Koenigs linearising function of the logistic map at
$w=\tfrac12$ and multiplier $r=2p$ (\cref{a:sec:koenigs}); and the dynamical
obstruction that explains the failure of the searches
(\cref{a:sec:ht,a:sec:scope}). Drawing on the classification of differentially
algebraic Schr\"oder functions due to Becker and Bergweiler
\cite{Becker1995}, \cref{a:sec:ht} proves that for every $r\in(0,1)$ the Koenigs
function $\psi_r$ is hypertranscendental in its own variable --- with no
exceptional parameters at all in the attracting window, since differentially
algebraic Schr\"oder functions exist only over repelling fixed points. The
classical elementary cases $r=2$ and $r=4$, both repelling and both outside the
branching range, are derived in \cref{a:app:integrable} and located inside the
classification.

The scope of the obstruction is stated with care in \cref{a:sec:scope}, because
it is easy to overstate. The theorem constrains the conjugacy as a function of its
dynamical variable $z$, and licenses no conclusion about individual values
$A(p_0)$ --- transcendence, and even irrationality, remain open --- nor about the
map $p\mapsto A(p)$, whose elementarity and differential algebraicity in $p$ are
likewise open. High-precision integer-relation searches at eleven rational
parameters, described in \cref{a:sec:scope}, supply the numerical evidence for the
conjecture that no closed form exists under any reading; they are finite evidence,
and are presented as such. \Cref{a:sec:practical} closes with the hybrid
product/asymptotic scheme used to evaluate $A(p)$ and $1/A(p)$ throughout the rest
of this thesis.
